The aim of this section is to prove the bounds on the degree of the primitive element $\theta$ of our number field, and the size of the constants used in the linear combination of the field generators that we use to construct it.

Let us first recall the setting for our number field. Given $k$ radicals $\mysqrt{d_1}{a_1},\ldots,\mysqrt{d_k}{a_k}$, with the $a_i$ pairwise coprime, and the $d_i$ and $a_i$ of magnitude at most $2^s$. Let $d=\lcm(d_1,\ldots,d_k)$, and denote by $\zeta_d$ a primitive $d$-th root of unity. We would like to construct the primitive element $\theta$ for the number field $K = \QQ(\mysqrt{d_1}{a_1},\ldots,\mysqrt{d_k}{a_k}, \zeta_d)$. Note that besides adjoining the radicals to $\QQ$, we also make sure to add $\zeta_d$, which ensures our field extension $K$ is Galois.

The primitive element for number fields is as follows:

\begin{theorem}\label{th:primitive_el_theorem_number_fields}

Let $K = L(\alpha_1,\ldots,\alpha_k)$ be a finite extension of $L$, and assume that $\alpha_2,\ldots,\alpha_k$ are separable over $L$. Then there is an element $\theta\in K$ such that $K=L(\theta)$.
\end{theorem}

The proof of the above theorem is constructive (see, e.g., \cite[Theorem~4.1.8]{cohen2013course} or \cite[Theorem 5.1]{milneGalTheory}), and computes the primitive element $\theta$ as a linear combination of the generators $\alpha_1,\ldots,\alpha_k$, that is $\theta = \sum_{i=1}^k c_i \alpha_i$. The computation of $\theta$ is to be done inductively, constructing first a primitive element $\theta_2$ for $L(\alpha_1,\alpha_2)$, then $\theta_3$ for $L(\alpha_1, \alpha_2,\alpha_3)$, and so on until $\theta_k$. Furthermore, it is shown that only finitely many combinations of the constants $c_i$ fail to generate a primitive element for the field extension $K$. This gives rise to an effective version of the primitive element theorem for number fields that induces a bound on the degree of the algebraic number $\theta$, as well as on the size of the constants $c_i$. In particular, for a field generated by two algebraic numbers, the bounds are as follow (see \cite[Proposition 6.6]{kururPhd}):

\begin{proposition}\label{prop:effective_primitive_el_th}
    Let $\alpha$ and $\beta$ be algebraic numbers of degree $m$ and $n$ respectively. There exists an integer $c\in\{1,\ldots,m^2n^2+1\}$ such that $\alpha+c\beta$ is a primitive element of $\QQ(\alpha,\beta)$.
\end{proposition}

Note that in general, we could choose $c_i\in\QQ$ with only finitely many combinations of the $c_i$'s not giving us a primitive element. Thus the primitive element need not be an algebraic integer in general. However, we make the choice to choose $c_i\in\ZZ$, therefore the minimal polynomial $f_\theta$ of $\theta$ is monic and $\theta$ an algebraic integer.

Now let us prove our claim on the bounds on our primitive element $\theta$:
\boundprimitiveelement*

\begin{proof}
Note that given $d_1,\ldots,d_k \leq 2^s$, their least common multiple is at most of size $2^{s^2}$, hence we have:
\[ \deg \zeta_d \leq  2^{s^2}.  \]

In the spirit of the proof of the primitive element theorem, use \Cref{prop:effective_primitive_el_th} inductively as follows:
\begin{align*}
    \theta_2 = \mysqrt{d_1}{a_1} + c_2\mysqrt{d_2}{a_2} &\qquad c_2 \leq (2^s)^2(2^s)^2 + 1 \textcolor{blue}{\leq 2^{4s}} \\
    \textcolor{blue}{\deg \theta_2 \leq 2^{2s}} \\
    \theta_3 = \mysqrt{d_1}{a_1} + c_2\mysqrt{d_2}{a_2} + c_3\mysqrt{d_3}{a_3}  &\qquad c_3 \leq (2^{2s})^2(2^s)^2 + 1 \textcolor{blue}{\leq 2^{6s}} \\
    \textcolor{blue}{\deg \theta_3 \leq 2^{3s}} \\
    \vdots & \\
    \theta_k = \mysqrt{d_1}{a_1} + c_2\mysqrt{d_2}{a_2} + \ldots + c_k\mysqrt{d_k}{a_k} &\qquad c_k \leq (2^{(k-1)s})^2(2^s)^2 + 1 \textcolor{blue}{\leq 2^{2s^2}}\\
    \textcolor{blue}{\deg \theta_k \leq 2^{s^2}} \\
    \theta = \theta_k + c_0 \zeta_d &\qquad c_0 \leq (2^{ks})^2(2^{s^2})^2 + 1 \textcolor{blue}{\leq 2^{4s^2}}\\
    \textcolor{blue}{\deg \theta_2 \leq 2^{2s^2}} \\
\end{align*}

The claimed bounds on the degree of $\theta$ and the size of the constants $c_i$ follow.
\end{proof}